% Chapter 8 — includable by main.tex
\chapter{Turingovi strojevi i odlučivost}
\label{chap:8}

\begin{quote}
\itshape
Sve što računala mogu raditi, mogu i Turingovi strojevi --- i obrnuto.
Ovo poglavlje definira taj univerzalni model izračuna, klasificira jezike
na odlučive i neodlučive, te dokazuje temeljne rezultate neodlučivosti
pomoću dijagonalizacije i svođenja.
\end{quote}


% ================================================================
\section{Turingov stroj --- motivacija}
% ================================================================

OA je mocniji od PA, ali jos uvijek ima ogranicenje: traka ne smije
biti dulja od ulaza. Što ako trebamo više prostora?

Najjednostavnija opcija: ``otvorimo'' traku s desne strane. Umjesto
da je traka duljine $|w|$, traka je \emph{beskonačna} prema desno,
inicijalno popunjena posebnim znakom praznine $\sqcup$.

Ovaj model zovemo \textbf{Turingov stroj} (TS), i ispostavlja se da je
to upravo dovoljno za modeliranje svakog izračuna koji je ikada
opisan. Ovo je sadržaj \textbf{Church-Turingove teze}: sve što je
``algoritamski'' izračunljivo, izračunljivo je i na Turingovom stroju.

TS dolazi u tri varijante s različitim svrhama:

\begin{center}
\renewcommand{\arraystretch}{1.4}
\begin{tabular}{lll}
\hline
\textbf{Vrsta} & \textbf{Svrha} & \textbf{Ključno} \\
\hline
Prepoznavac (TP) & prepoznaje jezik & može se vrtjeti unedogled \\
Odlučitelj (TO) & odlučuje problem & uvijek stane \\
Nabrajac (TE) & generira jezik & nema ulaza, nabrajamo izlaz \\
\hline
\end{tabular}
\end{center}

% ================================================================
\section{Turingov prepoznavac (TP)}
% ================================================================

\begin{definicija}[Turingov prepoznavac]
\textbf{Turingov prepoznavac} nad $\Sig$ je sedmorka
\[
  T = (Q,\, \Sig,\, \Gamma,\, \sqcup,\, \delta,\, q_0,\, q_X),
\]
gdje je $\Sig \subsetneq \Gamma$ (radna abeceda je sira od ulazne),
$\sqcup \in \Gamma \setminus \Sig$ je znak praznine, a
\[
  \delta : (Q \setminus \{q_X\}) \times \Gamma \to Q \times \Gamma \times \{-1, 1\}
\]
je \textbf{deterministicna} funkcija prijelaza.

Konfiguracija je trojka $(q, n, t) \in Q \times \mathbb{N} \times \Gamma^\mathbb{N}_\sqcup$,
gdje je $\Gamma^\mathbb{N}_\sqcup$ skup svih traka koje su od nekog mjesta
nadalje konstanta $\sqcup$ (prebrojivo mnogo takvih).

$T$ prihvaća riječ $w$ ako izračunavanje počevši od
$(q_0, 0, w\sqcup\sqcup\cdots)$ dosegne konfiguraciju sa stanjem $q_X$.
\end{definicija}

\begin{definicija}[Rekurzivno prebrojivi jezici]
Jezik je \textbf{rekurzivno prebrojiv} (u klasi RE) ako ga prepoznaje
neki TP. Piše se $L \in \mathrm{RE}$.
\end{definicija}

Ključna razlika od OA: TP se \emph{ne mora zaustaviti} za riječi koje
ne prihvaća. Može se vrtjeti unedogled. To je problem: ne znamo
je li stroj ``razmišlja'' ili je u beskonačnoj petlji.

% ================================================================
\section{Turingov odlučitelj (TO)}
% ================================================================

\begin{definicija}[Turingov odlučitelj i rekurzivni jezici]
\textbf{Turingov odlučitelj} (TO) je TP koji uz završno stanje $q_X$
ima i stanje odbijanja $\overline{q_X}$, te se \textbf{uvijek} zaustavi ---
ili u $q_X$ (prihvaćanje) ili u $\overline{q_X}$ (odbijanje).

Jezik je \textbf{rekurzivan} (u klasi Rek) ako ga prepoznaje neki TO.
\end{definicija}

Svaki TO je TP (možemo ignorirati $\overline{q_X}$), pa $\mathrm{Rek} \subseteq \mathrm{RE}$.

\begin{propozicija}[Postov teorem]
$L \in \mathrm{Rek}$ ako i samo ako su i $L$ i $L^c$ u $\mathrm{RE}$.
\end{propozicija}

\textit{Dokaz.}
$(\Rightarrow)$: ako je $L$ rekurzivan, i $L$ i $L^c$ su rekurzivni
(Rek je zatvorena na komplement), pa a fortiori u RE.

$(\Leftarrow)$: neka TP $S$ prepoznaje $L$, a TP $T$ prepoznaje $L^c$.
Konstruiramo TO koji istovremeno simulira $S$ i $T$, korak po korak.
Jer je svaka riječ ili u $L$ ili u $L^c$, jedan od njih će sigurno stati.
Kad $S$ prihvati, TO prihvaća; kad $T$ prihvati, TO odbija. \qed

\begin{kljucno}
Postov teorem znači: jezik je odlučiv točno onda kad možemo i
\emph{potvrditi} i \emph{zanijekati} članstvo algoritmom koji uvijek stane.
Samo znati "kako naći" nije dovoljno --- treba znati i "kako ne naći".
\end{kljucno}

% ================================================================
\section{Zatvorenost klasa Rek i RE}
% ================================================================

\begin{propozicija}
Klasa $\mathrm{Rek}$ zatvorena je na sve skupovne i jezicne operacije.
Klasa $\mathrm{RE}$ zatvorena je na $\cup$ i $\cap$, ali \textbf{ne} na komplement.
\end{propozicija}

Ključna tehnika za Rek: \textbf{potpuna simulacija}. Naredba
``simuliraj rad od $T$ na sadržaju trake $t_i$'' pokreće TO na
kopiji trake, te uvijek stane jer je $T$ odlučitelj.

Primjer --- zatvorenost Rek na Kleeneovu zvijezdu:
za riječ $w$, probamo sve moguće rastave $w = w_1 \cdots w_k$ (konačno
mnogo) i za svaki rastav provjerimo je li svaki $w_i \in L$. Ako ikoji
rastav uspije, prihvati; inače odbij.

Za RE s unijom: moramo simulirati oba TP \emph{naizmjenično}, korak
po korak --- ne smijemo pokrenuti prvi do kraja jer se možda nikad
ne zaustavi. Ovo je tzv.\\ "dovetailing" tehnika.

% ================================================================
\section{Kodiranje i problemi}
% ================================================================

Da bismo mogli govoriti o TS-ovima kao ulazima TS-u, trebamo ih
\textbf{kodirati} kao riječi. Koristimo \textbf{univerzalnu abecedu}
$\Sig = \{., \mathrm{I}, (, )\}$ i sustavno kodiramo:
\[
  c(n) := \mathrm{I}^n., \quad
  c(t_1, \ldots, t_k) := (c(t_1) \cdots c(t_k)), \quad
  c(\{t_1, \ldots, t_k\}) := (c(t_1) \cdots c(t_k)).
\]
Kod stroja $T$ (ili drugog matematickog objekta) oznacavamo $\langle T \rangle$,
a par $(T, w)$ oznacavamo $\langle T, w \rangle$.

\subsection{Prirodni strojevi i problemi}

\textbf{Problem} je jezik cija konkretna kodna abeceda nije bitna
\[
  A_\mathrm{KA} := \{\langle M, w \rangle : M \text{ je KA i } w \in L(M)\}
\]
je odlučiv problem (pišemo pseudokodom koji simulira KA $M$ na $w$).

% ================================================================
\section{Odlučivi i neodlučivi problemi}
% ================================================================

\subsection{Odlučivi problemi}

\begin{propozicija}
Sljedeći problemi su odlučivi:
\begin{itemize}
  \item $A_\mathrm{KA}$, $A_\mathrm{RI}$, $A_\mathrm{NKA}$ --- prihvaćanje za regularne jezike
  \item $E_\mathrm{KA}$ --- praznost KA-jezika
  \item $\mathrm{EQ}_\mathrm{KA}$ --- jednakost dva KA-jezika
  \item $A_\mathrm{BKG}$, $A_\mathrm{BK}$ --- prihvaćanje za BK-jezike (CYK!)
  \item $E_\mathrm{BKG}$, $E_\mathrm{BK}$ --- praznost BK-jezika
  \item $A_\mathrm{OA}$, $A_\mathrm{Kon}$ --- prihvaćanje za kontekstne jezike
\end{itemize}
\end{propozicija}

Svi ovi problemi imaju konkretne algoritme koje smo vec opisali
(pseudokodom ili formalnom konstrukcijom).

\subsection{Neodlučivi problemi --- hijerarhija}

\begin{center}
\renewcommand{\arraystretch}{1.5}
\begin{tabular}{lll}
\hline
\textbf{Problem} & \textbf{Klasa} & \textbf{Zašto} \\
\hline
$A_\mathrm{TP}$, $A_\mathrm{RE}$ & $\mathrm{RE} \setminus \mathrm{Rek}$ & TP može sebe simulirati \\
$\overline{A_\mathrm{TP}}$ & nije u RE & Postov teorem + $A_\mathrm{TP} \in \mathrm{RE} \setminus \mathrm{Rek}$ \\
\mathrm{EQ}_\mathrm{BK} & neodlučivo & svođenja od $\overline{A_\mathrm{TP}}$ \\
E_\mathrm{Kon} & neodlučivo & svođenja od $A_\mathrm{TP}$ \\
\hline
\end{tabular}
\end{center}

% ================================================================
\section{Svođenje (redukcija)}
% ================================================================

\begin{definicija}[Svođenje]
Kažemo da se $L$ \textbf{svodi na} $M$ ($L \leq M$) ako postoji
totalna Turing-izračunljiva funkcija $\varphi : \Sig^* \to \Sig^*$
takva da za svaki $w$:
\[
  w \in L \iff \varphi(w) \in M.
\]
\end{definicija}

\begin{propozicija}
Neka $L \leq M$.
\begin{itemize}
  \item Ako je $M$ (polu)odlučiv, tada je i $L$ (polu)odlučiv.
  \item Ako je $L$ neodlučiv, tada je i $M$ neodlučiv.
\end{itemize}
\end{propozicija}

Drugu točku koristimo za dokazivanje neodlučivosti:
Pokažemo $\mathrm{POZNATO\_NEODLUCIVO} \leq \mathrm{NOVI\_PROBLEM}$,
Čime zaključujemo da je i novi problem neodlučiv.

\begin{primjer}[Svođenje za $E_\mathrm{KA}$]
Problem jednakosti $\mathrm{EQ}_\mathrm{KA}$ svodi se na $E_\mathrm{KA}$:
\[
  \langle M_1, M_2 \rangle \mapsto \langle M_\Delta \rangle,
  \quad \text{gdje } L(M_\Delta) = L(M_1) \triangle L(M_2).
\]
Dakle $\langle M_1, M_2 \rangle \in \mathrm{EQ}_\mathrm{KA}
\iff \langle M_\Delta \rangle \in E_\mathrm{KA}$.
Konstruiranje $M_\Delta$ je mehaničko (Kartezijev produkt),
pa je $\varphi$ Turing-izračunljiva. Kako je $E_\mathrm{KA}$ odlučiv,
slijedi da je i $\mathrm{EQ}_\mathrm{KA}$ odlučiv. \checkmark
\end{primjer}

% ================================================================
\section{Dijagonalizacija: Russell i $A_\mathrm{TP}$}
% ================================================================

Kako dokazati da neki jezik \emph{nije} u RE? Lema o pumpanju
ne postoji za ovu klasu. Koristimo \textbf{dijagonalizaciju} --- trik
koji potjece od Cantora (prebrojiivost skupova) i Russella (paradoks).

\begin{propozicija}
Jezik
\[
  \mathrm{Russell} := \{\langle T \rangle : T \text{ je TP i } \langle T \rangle \notin L(T)\}
\]
nije rekurzivno prebrojiv.
\end{propozicija}

\textit{Dokaz.}
Pretpostavimo da je Russell u RE, i neka TP $R$ prepoznaje Russell.
Tada:
\[
  \langle R \rangle \in \mathrm{Russell}
  \iff \langle R \rangle \notin L(R) = \mathrm{Russell}.
\]
Kontradikcija. \qed

Ovo je direktan analog Russellovog paradoksa: ``skup svih skupova
koji ne sadrze sami sebe'' vodi do kontradikcije.

\begin{propozicija}
Jezik $A_\mathrm{TP} = \{\langle T, w \rangle : T \text{ prihvaća } w\}$
je u $\mathrm{RE} \setminus \mathrm{Rek}$.
\end{propozicija}

\textit{Skica dokaza.}
$A_\mathrm{TP} \in \mathrm{RE}$: konstruiramo \textbf{univerzalni Turingov stroj}
(UTS) koji simulira rad TP $T$ na ulazu $w$ (dan kao kod).
UTS je sam po sebi TP, pa je $A_\mathrm{TP} \in \mathrm{RE}$.

$A_\mathrm{TP} \notin \mathrm{Rek}$: dokazemo da $\overline{A_\mathrm{TP}}$ nije u RE.
Koristimo Russell $\leq \overline{A_\mathrm{TP}}$: svaki kod $\langle T \rangle$
preslikamo u $\langle T, \langle T \rangle \rangle$, i vrijedi:
\[
  \langle T \rangle \in \mathrm{Russell}
  \iff \langle T \rangle \notin L(T)
  \iff \langle T, \langle T \rangle \rangle \in \overline{A_\mathrm{TP}}.
\]
Dakle Russell $\leq \overline{A_\mathrm{TP}}$. Budući da Russell nije u RE,
ni $\overline{A_\mathrm{TP}}$ nije u RE. Po Postovom teoremu,
$A_\mathrm{TP} \notin \mathrm{Rek}$. \qed

% ================================================================
\section{Problem zaustavljanja i kompletna slika}
% ================================================================

\begin{propozicija}
Jezik $H = \{\langle T, w \rangle : \text{izračunavanje } T \text{ na } w \text{ stane}\}$
(``halting problem'') ekvivalentan je $A_\mathrm{TP}$ u smislu svedivosti:
$H \leq A_\mathrm{TP}$ i $A_\mathrm{TP} \leq H$.
\end{propozicija}

Ova ekvivalencija nije iznenadujuca: TP koji smo definirali stane
točno kada prihvate, pa "stajanje" i "prihvaćanje" su isti uvjet.

\subsection{Kompletna Chomskyjeva hijerarhija}

\begin{center}
\renewcommand{\arraystretch}{1.5}
\begin{tabular}{lllll}
\hline
\textbf{Klasa} & \textbf{Automat} & \textbf{Gram.} & \textbf{$A_L$} & \textbf{$E_L$} \\
\hline
Reg & KA & DLG & odlučiv & odlučiv \\
BK & PA & BKG & odlučiv & odlučiv \\
Kon & OA & KG/MG & odlučiv & neodlučiv \\
Rek & TO & --- & odlučiv & neodlučiv \\
RE & TP & OG & poluodlučiv & neodlučiv \\
\hline
\end{tabular}
\end{center}

Prave inkluzije potvrdjuju konkretni jezici:
\begin{align*}
  L_= &= \{a^n b^n\} \in \mathrm{BK} \setminus \mathrm{Reg} \\
  L_\equiv &= \{a^n b^n c^n\} \in \mathrm{Kon} \setminus \mathrm{BK} \\
  \mathrm{EQ}_\mathrm{RI_2} &\in \mathrm{Rek} \setminus \mathrm{Kon} \\
  A_\mathrm{TP} &\in \mathrm{RE} \setminus \mathrm{Rek} \\
  \overline{A_\mathrm{TP}} &\notin \mathrm{RE}
\end{align*}

% ================================================================
\section{Nedeterministicki Turingov stroj (NTP)}
% ================================================================

\begin{definicija}
NTP ima relaciju prijelaza umjesto funkcije. Prihvaća riječ
ako \emph{postoji barem jedna} grana izračunavanja koja stane.
\end{definicija}

\begin{propozicija}
NTP i TP su ekvivalentni: svaki NTP ima ekvivalentni TP.
\end{propozicija}

Dokaz: TP simulira NTP koristeci \textbf{tri trake}:
\begin{enumerate}
  \item Kopija ulaza (nepromjenjiva).
  \item Radna traka (simulacija jedne grane).
  \item Treća traka: kôd trenutne grane kao niz odabranih prijelaza
        iz $\Delta^*$ (leksikografskim redoslijedom isprobavamo sve grane).
\end{enumerate}
TP prolazi kroz sve moguce grane BFS-om: svaka iteracija produlji
granu za jedan korak. Čim neka grana prihvati, TP prihvaća.

\begin{napomena}
Za odlučitelje (NTO) analogna ekvivalencija vrijedi, ali je dokaz
teži --- ne može koristiti istu trotračnu simulaciju jer ne znamo
znači li $\overline{q_X}$ konačno odbijanje ili samo lokalni neuspjeh.
\end{napomena}

% ================================================================
\section{Turingovi nabrajaci (TE)}
% ================================================================

\begin{definicija}[Turingov nabrajac]
TE nema ulaza. Ima radnu traku i \textbf{izlaznu traku}, te posebno
\textbf{izlazno stanje} $q_\leftarrow$. Rad TE je uvijek beskonačni niz
konfiguracija. TE \textbf{izvodi} riječ $w$ ako ikad dosegne stanje
$q_\leftarrow$ s $w$ na izlaznoj traci.

$L(E) := \{w : E \text{ izvodi } w\}$.
\end{definicija}

\begin{propozicija}
$L \in \mathrm{RE}$ ako i samo ako postoji TE koji generira $L$.
\end{propozicija}

Smjer $(\Leftarrow)$: TP koji prepoznaje $L(E)$ simulira $E$ korak po
korak i uspoređuje izlaz s ulaznom riječi $w$. Čim se poklope,
prihvaća.

Smjer $(\Rightarrow)$: TE za $L$ (prepoznajan TP $T$) nabrajanjem:
za svaki $n$ i svaki moguc ulaz duljine $\leq n$, simulira $T$ na tom
ulazu točno $n$ koraka. Ako $T$ prihvati, izvede tu riječ.
Ovim ``dovetailing'' postupkom svaka riječ iz $L$ će biti izvedena
u konačno mnogo koraka.

% ================================================================
\section*{Sazetak}
% ================================================================

\medskip\noindent
\begin{tabular}{ll}
  TP & Turingov prepoznavač: stane ili prihvaća (nema garantije zaustavljanja) \\
  TO & Turingov odlučitelj: uvijek stane, prihvaća ili odbija \\
  RE & jezici koje prepoznaje neki TP \\
  Rek & jezici koje prepoznaje neki TO \\
  $A_\mathrm{TP}$ & poluodlučiv ali neodlučiv (svođenje + dijagonalizacija) \\
  Postov teorem & $L \in \mathrm{Rek} \iff L, L^c \in \mathrm{RE}$ \\
  Svođenje & $L \leq M$: odlučivost "prenosi se gore" \\
\end{tabular}

\medskip\noindent
\begin{center}
\Large\itshape
$\mathrm{Reg} \subsetneq \mathrm{BK} \subsetneq \mathrm{Kon} \subsetneq \mathrm{Rek} \subsetneq \mathrm{RE} \subsetneq \text{svi jezici}$
\end{center}